% Vietnamese translation for OI Public Library
% Source-File: IOI中国国家候选队论文/2004/周源.ppt
\documentclass[11pt]{article}
\usepackage{oipl-vn}

\newcommand{\Oh}{\mathcal{O}}

\title{Kết hợp số và hình trong thi tin học\\{\large Ghi chú theo slide}}
\author{Zhou Yuan}
\date{IOI 2004}

\begin{document}
\maketitle

\origtitle{Numeric-geometric thinking in informatics contests}
\sourcepath{IOI 2004 / Zhou Yuan.ppt}

\section{Tư tưởng chính}

Các slide nhấn mạnh hai chiều chuyển đổi: dùng hình để soi sáng quan hệ số,
và dùng số để xử lý đối tượng hình học. Khi một biểu thức đại số khó nhìn,
biểu diễn nó bằng điểm, đường thẳng hoặc độ dốc có thể làm lộ cấu trúc; khi
một bài hình học khó thao tác trực tiếp, tọa độ và thứ tự số lại giúp ta đưa
nó về bài toán tính toán.

\section{Dùng hình hỗ trợ số}

Mốc đầu tiên là bổ đề Raney: với dãy vòng có tổng bằng \(1\), đúng một cách
cắt làm mọi tổng tiền tố dương. Bài trình bày dùng hình ảnh đường gấp khúc
của tổng tiền tố để thấy vị trí cắt duy nhất.

Mốc tiếp theo là bài USACO về trung bình lớn nhất. Với tổng tiền tố
\(S_i\), giá trị
\[
  \operatorname{ave}(i,j)=\frac{S_j-S_{i-1}}{j-i+1}
\]
chính là hệ số góc của đoạn nối hai điểm
\((i-1,S_{i-1})\) và \((j,S_j)\). Điều kiện độ dài ít nhất \(F\) trở thành
ràng buộc khoảng cách theo trục hoành. Nhờ nhìn qua độ dốc, lời giải tránh
duyệt \(\Oh(n^2)\) trực tiếp và chuyển sang duy trì các ứng viên hình học.

\section{Dùng số hỗ trợ hình}

Các slide sau nhắc bài ``xưởng vẽ'' của POI: một mô hình điểm
\(P(x,y)\), phép tịnh tiến bởi \((X,Y)\), và thao tác trên dãy bit
\(a_1a_2a_3\), \(b_1b_2b_3\). Nội dung cốt lõi là mã hóa hình dạng bằng tọa
độ và trạng thái rời rạc, rồi dùng cấu trúc số để trả lời truy vấn nhanh.

\section{Ghi nhớ}

Trong cả hai hướng, bước quan trọng không phải là vẽ hình cho đẹp, mà là chọn
đại lượng hình học đúng: độ dốc, bao lồi, quan hệ nằm trên cùng một phía,
hoặc phép tịnh tiến. Khi đại lượng ấy khớp với công thức, phân tích độ phức
tạp mới trở nên rõ ràng.

\end{document}
